%==========================================================================
% Terceira Lista de Engenharia de Software RT
% Capítulos 22 e 23 — Ian Sommerville + Exercícios de Project Management
% (trabalho em grupo — 4 componentes)
%==========================================================================
\documentclass[12pt,a4paper]{article}

\usepackage[T1]{fontenc}
\usepackage[utf8]{inputenc}
\usepackage[brazil]{babel}
\usepackage{amsmath}
\usepackage{geometry}
\usepackage{xcolor}
\usepackage{tcolorbox}
\usepackage{enumitem}
\usepackage{tikz}
\usetikzlibrary{shapes.geometric,arrows.meta,positioning,calc,matrix}
\usepackage{booktabs}
\usepackage{array}
\usepackage{longtable}
\usepackage{multirow}
\usepackage{hyperref}

\geometry{top=2.2cm, bottom=2.2cm, left=2.8cm, right=2.8cm}
\setlength{\parskip}{5pt}
\setlength{\parindent}{0pt}

\tcbuselibrary{skins,breakable}
\newtcolorbox{questao}[1]{colback=gray!7,colframe=gray!45,
  fonttitle=\bfseries,title=#1,breakable}
\newtcolorbox{resposta}{colback=blue!4,colframe=blue!50!black,
  fonttitle=\bfseries,title=Resposta,breakable}

%==========================================================================
\begin{document}

\begin{center}
  {\Large\bfseries Terceira Lista de Engenharia de Software RT}\\[5pt]
  {\large Capítulos 22 e 23 — Gerenciamento de Projetos de Software}\\[4pt]
  {\normalsize Ian Sommerville — \textit{Software Engineering}, 10ª ed.
   \quad|\quad Trabalho em grupo (4 componentes)}
\end{center}
\hrule\vspace{6pt}

%==========================================================================
\section*{Capítulo 22 — Gerenciamento de Projetos}
%==========================================================================

%--------------------------------------------------------------------------
\subsection*{Questão 22.1}
%--------------------------------------------------------------------------
\begin{questao}{Enunciado}
Explique por que a intangibilidade dos sistemas de software gera problemas
especiais para gerenciamento de projetos de software.
\end{questao}
\begin{resposta}
Ao contrário de sistemas físicos, o software não é visível: não é possível
inspecionar um executável e avaliar seu progresso ou qualidade a olho nu.
Isso cria três classes de problemas gerenciais:

\begin{enumerate}[nosep]
  \item \textbf{Dificuldade de medir progresso:} marcos tangíveis (como um
        componente mecânico entregue) não existem. Gerentes precisam confiar
        em métricas indiretas (linhas de código, pontos de função, testes
        passados), que são facilmente distorcidas. A síndrome ``90\%
        completo'' é comum — o produto parece quase pronto por semanas.

  \item \textbf{Invisibilidade de defeitos:} erros de lógica, condições de
        corrida e vulnerabilidades de segurança não se manifestam visualmente,
        podendo permanecer latentes até situações específicas de operação.
        Isso torna difícil garantir qualidade ao longo do processo.

  \item \textbf{Gerenciamento de contratos e expectativas:} clientes não
        conseguem avaliar o que foi entregue sem executar o sistema. Isso
        complica a aceitação formal, o pagamento por milestones e a
        negociação de mudanças de escopo.

  \item \textbf{Estimativas imprecisas:} sem componentes físicos de referência,
        estimar custo e prazo depende de experiência anterior com sistemas
        similares — dados raros em domínios novos.
\end{enumerate}
\end{resposta}

%--------------------------------------------------------------------------
\subsection*{Questão 22.2}
%--------------------------------------------------------------------------
\begin{questao}{Enunciado}
Explique por que os melhores programadores nem sempre se tornam os melhores
gerentes de software.
\end{questao}
\begin{resposta}
Gerenciamento e programação são conjuntos de habilidades ortogonais.
Conforme a Seção~22.1 de Sommerville, as principais atividades de um gerente
incluem: escrever propostas, planejamento, monitoração, relatórios e
gerenciamento de pessoas. Programação exige raciocínio lógico profundo,
concentração em problemas técnicos e trabalho predominantemente individual.

Conflitos específicos:
\begin{itemize}[nosep]
  \item \textbf{Comunicação interpessoal:} gerentes precisam negociar, motivar
        equipes e transmitir informações a públicos não técnicos. Bons
        programadores frequentemente preferem trabalho solitário e comunicação
        direta e técnica.
  \item \textbf{Tolerância à ambiguidade:} problemas gerenciais raramente têm
        uma solução única e correta. Programadores costumam ser treinados a
        buscar soluções ótimas, o que pode paralisar decisões gerenciais
        necessariamente baseadas em informação incompleta.
  \item \textbf{Perda de contribuição técnica:} gerentes deixam de codificar.
        Programadores excepcionais podem perceber a promoção como perda de
        produtividade e satisfação pessoal.
  \item \textbf{Responsabilidade difusa:} o sucesso de um time não é
        atribuível a um único indivíduo, o que contrasta com o reconhecimento
        direto que bons programadores recebem por soluções técnicas elegantes.
\end{itemize}
\end{resposta}

%--------------------------------------------------------------------------
\subsection*{Questão 22.3}
%--------------------------------------------------------------------------
\begin{questao}{Enunciado}
Usando instâncias de problemas de projetos relatados na literatura, liste as
dificuldades e os erros de gerenciamento que ocorreram em projetos de
programação que falharam.
\end{questao}
\begin{resposta}
Com base em \textit{The Mythical Man-Month} (Brooks, 1975) e outros casos:

\begin{itemize}
  \item \textbf{Lei de Brooks (IBM OS/360):} adicionar pessoal a um projeto
        atrasado atrasa-o ainda mais. O overhead de comunicação cresce em
        $O(n^2)$ com o número de pessoas. A IBM estimou errado a complexidade
        do OS/360 e, ao adicionar equipes, multiplicou os atrasos.

  \item \textbf{Estimativa otimista (segundo sistema de Brooks):} o segundo
        sistema de um desenvolvedor tende ao excesso de complexidade e
        funcionalidades. O OS/360 sofreu disso — funcionalidades desnecessárias
        foram adicionadas, aumentando a complexidade sem valor correspondente.

  \item \textbf{Denver Airport Baggage System (1994):} \$193 milhões em
        prejuízo e 16 meses de atraso por subestimação da complexidade de
        integração, falta de prototipação e pressões políticas sobre o
        cronograma.

  \item \textbf{Healthcare.gov (2013):} falha no lançamento do portal do
        governo americano por gestão fragmentada entre múltiplos
        contratantes, falta de um arquiteto-chefe e testes insuficientes
        de integração até semanas antes do lançamento.

  \item \textbf{Erros recorrentes identificados:} cronogramas não realistas
        sob pressão comercial; ausência de controle de mudanças formal;
        falta de comunicação entre subequipes; testes relegados ao final
        do projeto; mitos de produtividade (``pessoas-mês'' como unidade
        de medida de capacidade).
\end{itemize}
\end{resposta}

%--------------------------------------------------------------------------
\subsection*{Questão 22.4}
%--------------------------------------------------------------------------
\begin{questao}{Enunciado}
Além dos riscos mostrados na Tabela~22.1, identifique pelo menos seis outros
possíveis riscos que podem surgir em projetos de software.
\end{questao}
\begin{resposta}
\begin{enumerate}[nosep]
  \item \textbf{Risco de segurança cibernética:} vulnerabilidades descobertas
        em componentes de terceiros (ex.: Log4Shell) podem invalidar decisões
        arquiteturais tardias no projeto.
  \item \textbf{Risco de integração de sistemas legados:} interfaces não
        documentadas ou indisponíveis de sistemas legados podem impedir a
        migração ou integração planejada.
  \item \textbf{Risco de rotatividade de pessoal-chave:} saída de membros
        com conhecimento crítico não documentado pode paralisar o projeto.
  \item \textbf{Risco regulatório:} mudanças em legislação (ex.: LGPD,
        GDPR) durante o desenvolvimento podem exigir retrabalho arquitetural
        significativo.
  \item \textbf{Risco de fornecedor/terceiro:} falência ou mudança de
        estratégia de um fornecedor de componente crítico (ex.: descontinuação
        de API) pode inviabilizar partes do sistema.
  \item \textbf{Risco de desempenho:} o sistema atende aos requisitos
        funcionais mas não às metas de desempenho sob carga real, exigindo
        refatoração extensiva.
  \item \textbf{Risco de adoção:} usuários resistem ao novo sistema por
        questões de usabilidade ou mudança de processos, reduzindo o retorno
        sobre o investimento.
\end{enumerate}
\end{resposta}

%--------------------------------------------------------------------------
\subsection*{Questão 22.5}
%--------------------------------------------------------------------------
\begin{questao}{Enunciado}
Contratos de preço fixo podem mover o risco de projeto para o contratante.
Sugira como o uso desses contratos pode aumentar a probabilidade de surgirem
riscos de produtos.
\end{questao}
\begin{resposta}
Em contratos de preço fixo, o contratante (desenvolvedor) absorve qualquer
custo acima do valor contratado. Para proteger sua margem, o contratante
tende a:

\begin{itemize}[nosep]
  \item \textbf{Cortar escopo de qualidade:} reduzir testes, revisões de
        código e documentação — atividades que não são visíveis ao cliente
        mas são críticas para a confiabilidade do produto final.
  \item \textbf{Aceitar requisitos ambíguos sem negociá-los:} cada
        esclarecimento posterior pode ser cobrado como mudança, incentivando
        a entrega de um produto que atende à letra, mas não ao espírito dos
        requisitos.
  \item \textbf{Usar tecnologias ou arquiteturas de menor custo de
        desenvolvimento:} que podem não ser as mais adequadas para
        manutenibilidade ou desempenho a longo prazo.
  \item \textbf{Subestimar deliberadamente:} para ganhar o contrato,
        contratantes podem oferecer lances inviáveis, sabendo que a
        recuperação virá por meio de cobranças de mudanças ou entregando
        um produto de qualidade inferior.
\end{itemize}

Todos esses comportamentos aumentam o risco de produto: o sistema entregue
pode ser frágil, inseguro ou não atender às necessidades reais, mesmo
cumprindo formalmente o contrato.
\end{resposta}

%--------------------------------------------------------------------------
\subsection*{Questão 22.6}
%--------------------------------------------------------------------------
\begin{questao}{Enunciado}
Explique por que manter todos os membros de um grupo informados sobre
progressos e decisões técnicas de um projeto pode melhorar a coesão de grupo.
\end{questao}
\begin{resposta}
A coesão de grupo é fortalecida quando todos os membros compartilham um
modelo mental comum do projeto. Informar progressos e decisões técnicas:

\begin{enumerate}[nosep]
  \item \textbf{Gera senso de pertencimento:} membros informados sentem que
        contribuem para um objetivo compartilhado, não apenas para sua tarefa
        isolada. Isso reduz o comportamento de ``silo''.
  \item \textbf{Cria responsabilidade coletiva:} quando todos sabem o estado
        do projeto, erros e atrasos são detectados mais cedo e a equipe age
        de forma proativa.
  \item \textbf{Facilita decisões distribuídas:} membros informados podem
        tomar micro-decisões alinhadas às decisões estratégicas sem precisar
        escalar cada questão.
  \item \textbf{Reduz rumores e ansiedade:} transparência sobre dificuldades
        e riscos evita o ``jogo de telefone'' que frequentemente causa
        conflitos e desmotivação.
  \item \textbf{Promove aprendizado coletivo:} decisões técnicas explicadas
        à equipe aumentam a competência coletiva e reduzem a dependência
        de especialistas individuais.
\end{enumerate}
\end{resposta}

%--------------------------------------------------------------------------
\subsection*{Questão 22.7}
%--------------------------------------------------------------------------
\begin{questao}{Enunciado}
Quais os problemas que podem surgir nas equipes de Extreme Programming em
que muitas decisões de gerenciamento são atribuídas aos membros de equipe?
\end{questao}
\begin{resposta}
\begin{itemize}[nosep]
  \item \textbf{Conflitos não resolvidos:} sem uma autoridade clara,
        divergências técnicas (ex.: escolha de arquitetura, prioridade de
        histórias) podem se arrastar sem decisão, criando bloqueios.
  \item \textbf{Planejamento inconsistente:} diferentes membros podem
        ter visões conflitantes sobre velocidade de desenvolvimento e
        estimativas de story points, dificultando o planejamento de releases.
  \item \textbf{Falta de accountability:} quando decisões são coletivas,
        ninguém é responsável pelos resultados. Erros gerenciais podem não
        ser reconhecidos ou corrigidos.
  \item \textbf{Pressão grupal:} membros menos assertivos podem silenciar
        diante de colegas mais dominantes, perdendo-se perspectivas valiosas.
  \item \textbf{Dificuldade com stakeholders externos:} clientes e gestores
        precisam de um ponto de contato único. Uma equipe sem hierarquia
        clara pode comunicar mensagens inconsistentes.
  \item \textbf{Escalabilidade:} o modelo funciona bem em equipes pequenas
        (até 10 pessoas). Em projetos maiores, a autogestão coletiva torna-se
        caótica sem mecanismos formais de coordenação.
\end{itemize}
\end{resposta}

%--------------------------------------------------------------------------
\subsection*{Questão 22.8}
%--------------------------------------------------------------------------
\begin{questao}{Enunciado}
Escreva um estudo de caso, no estilo usado neste capítulo, para ilustrar a
importância da comunicação em uma equipe de projeto com membros remotos.
\end{questao}
\begin{resposta}
\textbf{Caso: Sistema de Monitoramento de Infraestrutura Hospitalar}

Uma equipe de 8 desenvolvedores distribuídos em Manaus, São Paulo e Lisboa
foi contratada para desenvolver um sistema de monitoramento de leitos
hospitalares em 9 meses. O time de Manaus era responsável pelo backend
(sensores IoT), o time de São Paulo pelo frontend web e o time de Lisboa
pela integração com o sistema de prontuários.

\textbf{Problema:} Os times de Manaus e Lisboa definiram, independentemente,
formatos de mensagem JSON diferentes para os eventos de sensor. A discrepância
só foi descoberta na integração (mês 7), quando o frontend exibia dados
incorretos. A correção exigiu 3 semanas adicionais e refatoração do
módulo de persistência.

\textbf{Causa raiz:} As reuniões semanais de sincronização focavam em
progresso funcional individual, sem revisão de contratos de API. Não havia
um repositório central de especificações de interface revisado por todos
os times. As diferenças de fuso horário (Manaus--Lisboa: -4h) dificultavam
comunicação assíncrona e resolução rápida de dúvidas.

\textbf{Lição:} Equipes distribuídas necessitam de um contrato de interface
explícito (ex.: OpenAPI), revisado e aprovado por todos os times antes do
início da implementação. Reuniões de sincronização devem incluir revisão
de decisões técnicas de impacto transversal, não apenas relatórios de
progresso. Uma ferramenta de comunicação com canais por componente e
responsável único por decisões técnicas teria prevenido o problema.
\end{resposta}

%--------------------------------------------------------------------------
\subsection*{Questão 22.9}
%--------------------------------------------------------------------------
\begin{questao}{Enunciado}
Seu gerente solicitou que você entregue o software dentro do cronograma,
o qual só poderá ser cumprido caso você peça à sua equipe que trabalhe horas
extras não remuneradas. Todos os membros têm crianças pequenas. Discuta se
você deve aceitar essa demanda ou convencer sua equipe a dar seu tempo à
organização.
\end{questao}
\begin{resposta}
\textbf{Fatores a considerar:}

\begin{itemize}[nosep]
  \item \textbf{Ético-profissional:} o Código de Ética da ACM e da IEEE-CS
        exige que engenheiros tratem colegas com justiça e não permitam que
        pressões comerciais comprometam o bem-estar da equipe.
  \item \textbf{Responsabilidade parental:} solicitar horas extras a membros
        com filhos pequenos impõe custo social real às suas famílias, que não
        têm voz na negociação contratual.
  \item \textbf{Qualidade do software:} equipes exaustas cometem mais erros.
        Horas extras não remuneradas frequentemente resultam em dívida técnica
        que, a longo prazo, atrasa projetos futuros.
  \item \textbf{Comprometimento contratual:} se o cronograma foi estabelecido
        com estimativas ruins, o problema é de planejamento — não de
        esforço da equipe.
\end{itemize}

\textbf{Decisão recomendada:} Não aceitar a demanda silenciosamente. A postura
ética é comunicar ao gerente o impacto humano e técnico da solicitação,
propor renegociação do escopo ou do prazo com o cliente, e documentar
formalmente as consequências da decisão. Se a demanda for imposta,
comunicá-la à equipe de forma transparente — sem coerção velada — e deixar
que cada membro decida individualmente se pode ou não contribuir com tempo
adicional voluntário. Jamais pressionar quem recusar.
\end{resposta}

%--------------------------------------------------------------------------
\subsection*{Questão 22.10}
%--------------------------------------------------------------------------
\begin{questao}{Enunciado}
Como programador, você tem a possibilidade de ser promovido para gerente de
projetos, mas sente que pode contribuir mais eficazmente como técnico. Discuta
se você deve aceitar a promoção.
\end{questao}
\begin{resposta}
\textbf{Argumentos para recusar:}
\begin{itemize}[nosep]
  \item Contribuição técnica de alta qualidade tem impacto direto e mensurável
        em sistemas críticos; gerenciamento ruim pode neutralizar esse impacto.
  \item Satisfação pessoal e motivação intrínseca são preditores fortes de
        produtividade a longo prazo. Assumir um papel que não mobiliza as
        competências valorizadas pelo profissional tende a gerar frustração.
  \item Muitas organizações precisam mais de engenheiros sêniores do que
        de gerentes adicionais.
\end{itemize}

\textbf{Argumentos para aceitar:}
\begin{itemize}[nosep]
  \item Gerentes com sólida base técnica tomam decisões melhores e ganham
        credibilidade da equipe. A organização pode se beneficiar mais de
        alguém que entende os detalhes técnicos do que de um gerente puramente
        administrativo.
  \item Exposição gerencial pode revelar aptidões não descobertas e abrir
        novas perspectivas de carreira.
\end{itemize}

\textbf{Recomendação:} negociar com a organização alternativas como papel de
\textit{Tech Lead} — responsabilidade técnica e alguma responsabilidade de
coordenação, sem abandono da contribuição técnica. Se a promoção for aceita,
deve-se fazê-lo com comprometimento genuíno e plano de desenvolvimento
das habilidades gerenciais. Aceitar a promoção sem motivação apenas para
progressão salarial é prejudicial ao profissional, à equipe e ao projeto.
\end{resposta}

%==========================================================================
\section*{Capítulo 23 — Planejamento de Projetos}
%==========================================================================

%--------------------------------------------------------------------------
\subsection*{Questão 23.1}
%--------------------------------------------------------------------------
\begin{questao}{Enunciado}
Em quais circunstâncias uma empresa poderia, justificadamente, cobrar um
preço muito mais elevado do que o custo estimado de software, além de uma
margem de lucro razoável?
\end{questao}
\begin{resposta}
\begin{enumerate}[nosep]
  \item \textbf{Monopólio de conhecimento específico:} se a empresa é a única
        com expertise no domínio do cliente (ex.: sistema de defesa ou
        industrial específico), pode precificar o valor percebido, não o custo.
  \item \textbf{Risco assumido pelo fornecedor:} em contratos de preço fixo
        com requisitos altamente incertos, o prêmio de risco justifica uma
        margem maior.
  \item \textbf{Software crítico de segurança:} sistemas médicos, aviônicos
        e nucleares exigem certificação formal, documentação extensiva e
        processos de garantia de qualidade caros que elevam os custos reais.
  \item \textbf{Entrega urgente:} prazos muito comprimidos exigem horas extras,
        contratação rápida ou interrupção de outros projetos, justificando
        preço premium.
  \item \textbf{Propriedade intelectual:} quando o software incorpora
        algoritmos proprietários ou pesquisa avançada, o valor intelectual
        supera o custo de desenvolvimento.
\end{enumerate}
\end{resposta}

%--------------------------------------------------------------------------
\subsection*{Questão 23.2}
%--------------------------------------------------------------------------
\begin{questao}{Enunciado}
Explique por que o processo de planejamento de projeto é iterativo e por que
um plano deve ser continuamente revisto durante um projeto de software.
\end{questao}
\begin{resposta}
O planejamento é iterativo porque as premissas iniciais são inevitavelmente
incompletas ou incorretas:

\begin{itemize}[nosep]
  \item \textbf{Incerteza inicial:} no início do projeto, requisitos são
        vagos, tecnologias não testadas e a capacidade da equipe é incerta.
        Estimativas baseadas nessas premissas degradam-se rapidamente.
  \item \textbf{Descobertas durante o desenvolvimento:} cada fase revela
        complexidades não previstas. Um módulo estimado em 5 dias pode levar
        15 após análise detalhada.
  \item \textbf{Mudanças externas:} requisitos do cliente mudam; membros de
        equipe entram ou saem; tecnologias de terceiros mudam de versão.
  \item \textbf{Feedback de milestones:} revisões periódicas revelam desvios
        entre o planejado e o executado, exigindo replanejamento.
\end{itemize}

Um plano não revisado torna-se desatualizado e perde seu valor como
instrumento de controle. O plano iterativo é mais eficaz porque incorpora
aprendizado contínuo e mantém as estimativas ancoradas na realidade corrente
do projeto.
\end{resposta}

%--------------------------------------------------------------------------
\subsection*{Questão 23.3}
%--------------------------------------------------------------------------
\begin{questao}{Enunciado}
Explique rapidamente a finalidade de cada uma das seções em um plano de
projeto de software.
\end{questao}
\begin{resposta}
\begin{center}
\begin{tabular}{p{4cm}p{9.5cm}}
\toprule
\textbf{Seção} & \textbf{Finalidade} \\
\midrule
Introdução & Define escopo, objetivos e contexto; estabelece quem vai ler o
             documento e o que se espera do projeto. \\
\midrule
Organização do projeto & Descreve a estrutura da equipe, papéis e responsabilidades,
                         e a relação com o cliente. \\
\midrule
Análise de riscos & Identifica riscos, sua probabilidade e impacto, e define
                    estratégias de mitigação e contingência. \\
\midrule
Req. de recursos de hardware e software & Lista ferramentas, infraestrutura e
                                          ambientes necessários para desenvolvimento
                                          e operação. \\
\midrule
Divisão do trabalho & Decompõe o projeto em atividades, atribui responsáveis
                      e define dependências (WBS). \\
\midrule
Cronograma do projeto & Presenta o sequenciamento temporal das atividades (Gantt,
                        CPM), marcos e datas de entrega. \\
\midrule
Mecanismos de monitoramento e relatórios & Define como o progresso será medido,
                                           com que frequência e em que formato
                                           relatórios serão produzidos. \\
\bottomrule
\end{tabular}
\end{center}
\end{resposta}

%--------------------------------------------------------------------------
\subsection*{Questão 23.4}
%--------------------------------------------------------------------------
\begin{questao}{Enunciado}
Sugira quatro maneiras de reduzir o risco em uma estimativa de custo.
\end{questao}
\begin{resposta}
\begin{enumerate}[nosep]
  \item \textbf{Usar múltiplas técnicas de estimativa:} combinar COCOMO,
        pontos de função e julgamento especializado. Convergência entre os
        métodos aumenta a confiança; divergência sinaliza incerteza que
        deve ser explicitada como risco.
  \item \textbf{Estimativas incrementais (top-down + bottom-up):} decompor o
        sistema em componentes e estimar cada um individualmente (bottom-up),
        comparando com uma estimativa global (top-down). Discrepâncias revelam
        itens omitidos.
  \item \textbf{Usar dados históricos de projetos similares:} repositórios
        organizacionais de projetos passados (com tamanho, custo real e
        desvios) fornecem base empírica para calibrar estimativas.
  \item \textbf{Inclusão explícita de reservas de contingência:} reservas
        para risco (ex.: 20--30\% do esforço estimado) transformam a
        incerteza de uma ameaça implícita em um buffer gerenciado e visível
        para o cliente.
\end{enumerate}
\end{resposta}

%--------------------------------------------------------------------------
\subsection*{Questão 23.5 — Gráfico de Barras da Tabela 23.7}
%--------------------------------------------------------------------------
\begin{questao}{Enunciado}
A Tabela~23.7 define tarefas com durações e dependências. Desenhe um gráfico
de barras mostrando o cronograma de projeto.
\end{questao}
\begin{resposta}
Cronograma calculado (início mais cedo para cada tarefa):

\begin{center}\small
\begin{tabular}{cccp{5cm}}
\toprule
\textbf{Tarefa} & \textbf{ES} & \textbf{EF} & \textbf{Dependências} \\
\midrule
T1 & 0 & 10 & --- \\
T2 & 10 & 25 & T1 \\
T3 & 25 & 35 & T1, T2 \\
T4 & 0 & 20 & --- \\
T5 & 0 & 10 & --- \\
T6 & 35 & 50 & T3, T4 \\
T7 & 35 & 55 & T3 \\
T8 & 55 & 90 & T7 \\
T9 & 50 & 65 & T6 \\
T10 & 65 & 70 & T5, T9 \\
T11 & 65 & 75 & T9 \\
T12 & 70 & 90 & T10 \\
T13 & 35 & 70 & T3, T4 \\
T14 & 90 & 100 & T8, T9 \\
T15 & 100 & 120 & T2, T14 \\
T16 & 120 & 130 & T15 \\
\bottomrule
\end{tabular}
\quad \textbf{Duração total: 130 dias}
\end{center}

\textbf{Gráfico de Gantt (escala: 1 unidade = 10 dias):}

\begin{center}
\begin{tikzpicture}[xscale=0.09, yscale=0.38]
  \foreach \x in {0,10,...,130} {
    \draw[gray!30,thin] (\x,-0.3) -- (\x,17.3);
    \node[below,font=\tiny] at (\x,-0.3) {\x};
  }
  \foreach \row/\name/\s/\f in {
      1/T1/0/10, 2/T2/10/25, 3/T3/25/35, 4/T4/0/20,
      5/T5/0/10, 6/T6/35/50, 7/T7/35/55, 8/T8/55/90,
      9/T9/50/65, 10/T10/65/70, 11/T11/65/75, 12/T12/70/90,
      13/T13/35/70, 14/T14/90/100, 15/T15/100/120, 16/T16/120/130} {
    \pgfmathsetmacro{\y}{17-\row}
    \fill[blue!50!cyan,draw=blue!70,rounded corners=1pt]
      (\s,\y+0.1) rectangle (\f,\y+0.85);
    \node[left,font=\tiny\bfseries] at (0,\y+0.48) {\name};
  }
  \draw[->] (0,-0.3) -- (135,-0.3) node[right,font=\small] {Dias};
  \draw (0,-0.3) -- (0,17.3);
  \node[above,font=\small\bfseries] at (65,17.5) {Cronograma — Tabela 23.7};
\end{tikzpicture}
\end{center}
\end{resposta}

%--------------------------------------------------------------------------
\subsection*{Questão 23.6 — Reorganização com T5 = 40 dias}
%--------------------------------------------------------------------------
\begin{questao}{Enunciado}
Suponha que ocorra um retrocesso grave e, em vez de levar 10 dias, a tarefa
T5 leva 40 dias. Elabore novos gráficos de barras, mostrando como o projeto
pode ser reorganizado.
\end{questao}
\begin{resposta}
Com T5 = 40 dias:
\begin{itemize}[nosep]
  \item T5: EF passa de 10 para \textbf{40 dias}.
  \item T10 depende de T5 e T9. Com T9 terminando no dia 65, a restrição
        de T10 continua sendo T9:
        $\text{ES}(\text{T10}) = \max(40, 65) = 65$.
  \item \textbf{Nenhuma outra tarefa é afetada} — T5 não é caminho crítico.
  \item \textbf{Duração total do projeto permanece 130 dias.}
\end{itemize}

\textbf{Impacto:} apenas o comprimento da barra de T5 no Gantt aumenta.
A reorganização consiste em garantir que os recursos alocados a T5 sejam
mantidos por 40 dias em vez de 10, sem impacto nas demais atividades.

\begin{center}
\begin{tikzpicture}[xscale=0.09, yscale=0.38]
  \foreach \x in {0,10,...,130} {
    \draw[gray!30,thin] (\x,-0.3) -- (\x,17.3);
    \node[below,font=\tiny] at (\x,-0.3) {\x};
  }
  \foreach \row/\name/\s/\f in {
      1/T1/0/10, 2/T2/10/25, 3/T3/25/35, 4/T4/0/20,
      5/T5/0/10, 6/T6/35/50, 7/T7/35/55, 8/T8/55/90,
      9/T9/50/65, 10/T10/65/70, 11/T11/65/75, 12/T12/70/90,
      13/T13/35/70, 14/T14/90/100, 15/T15/100/120, 16/T16/120/130} {
    \pgfmathsetmacro{\y}{17-\row}
    \fill[blue!50!cyan,draw=blue!70,rounded corners=1pt]
      (\s,\y+0.1) rectangle (\f,\y+0.85);
    \node[left,font=\tiny\bfseries] at (0,\y+0.48) {\name};
  }
  % T5 extended to 40d (row 5 -> y=12)
  \fill[red!60,draw=red!80,rounded corners=1pt]
    (0,12.1) rectangle (40,12.85);
  \node[font=\tiny\bfseries,white] at (20,12.48) {T5 (40d)};
  \draw[->] (0,-0.3) -- (135,-0.3) node[right,font=\small] {Dias};
  \draw (0,-0.3) -- (0,17.3);
  \node[above,font=\small\bfseries] at (65,17.5) {Cronograma Reorganizado (T5=40 dias)};
\end{tikzpicture}
\end{center}
\end{resposta}

%--------------------------------------------------------------------------
\subsection*{Questão 23.7}
%--------------------------------------------------------------------------
\begin{questao}{Enunciado}
Explique os potenciais problemas com o jogo de planejamento do XP quando o
software tem altos requisitos de desempenho ou de confiança.
\end{questao}
\begin{resposta}
O jogo de planejamento do XP baseia-se em histórias de usuário priorizadas
pelo cliente por valor de negócio. Problemas em sistemas de alto desempenho
ou alta confiança:

\begin{itemize}[nosep]
  \item \textbf{Requisitos não-funcionais invisíveis nas histórias:} desempenho,
        segurança e confiabilidade raramente aparecem como histórias de usuário
        discretas. São propriedades emergentes da arquitetura e difíceis de
        priorizar por valor de negócio percebido.
  \item \textbf{Decisões arquiteturais adiadas:} o XP encoraja o design emergente
        (YAGNI), mas sistemas de alta confiança precisam de uma arquitetura
        de segurança definida cedo — mudar tarde é extremamente custoso.
  \item \textbf{Ausência de análise formal de riscos:} histórias de segurança
        crítica (ex.: ``o sistema não deve emitir dose de radiação errada'')
        não são priorizáveis apenas por valor de negócio; são obrigações legais.
  \item \textbf{Ciclos curtos vs. certificação:} certificações de sistemas
        críticos (DO-178C, IEC~61508) exigem documentação de processo pesada,
        incompatível com os ciclos semanais e baixa cerimônia do XP.
  \item \textbf{Refatoração não é livre:} em sistemas de alto desempenho,
        refatorar uma estrutura de dados pode quebrar garantias de tempo de
        resposta garantidas previamente.
\end{itemize}
\end{resposta}

%--------------------------------------------------------------------------
\subsection*{Questão 23.8}
%--------------------------------------------------------------------------
\begin{questao}{Enunciado}
Sistema de radioterapia com estimativa COCOMO de 26 pessoas-mês, com todos
os multiplicadores de driver de custo definidos como 1. Explique por que
essa estimativa deve ser ajustada, cite quatro fatores e proponha valores.
\end{questao}
\begin{resposta}
Definir todos os multiplicadores como 1 assume que o sistema é \emph{médio}
em todas as dimensões — o que é claramente inadequado para um sistema médico
crítico embarcado. Os fatores COCOMO II que precisam de ajuste:

\begin{center}
\small
\begin{tabular}{p{3.5cm}p{5.5cm}cc}
\toprule
\textbf{Fator} & \textbf{Justificativa} & \textbf{Valor} & \textbf{Multiplicador} \\
\midrule
RELY (Confiabilidade) & Sistema de radioterapia — falha pode resultar em
                        morte do paciente. Nível ``extra alto''. & Extra alto & 1,39 \\
\midrule
CPLX (Complexidade) & Controle em tempo real de hardware específico
                      (processador dedicado, 256 MB RAM fixos), comunicação
                      com BD hospitalar, interface crítica de segurança.
                      Nível ``muito alto''. & Muito alto & 1,30 \\
\midrule
STOR (Restrição de armazenamento) & Execução em memória fixa de 256 MB com
                      código de controle em tempo real. Restrição ``muito
                      alta'' (uso acima de 70\% da memória). & Muito alto & 1,21 \\
\midrule
SECU (Segurança) & Restrição de segurança funcional (requisito de integridade
                   de segurança nível SIL 3 ou 4). Exige análise formal,
                   FMEA e rastreabilidade completa. & Extra alto & 1,30 \\
\bottomrule
\end{tabular}
\end{center}

Produto dos multiplicadores: $1{,}39 \times 1{,}30 \times 1{,}21 \times 1{,}30 \approx 2{,}84$

\textbf{Estimativa ajustada:} $26 \times 2{,}84 \approx \mathbf{74}$ pessoas-mês.
\end{resposta}

%--------------------------------------------------------------------------
\subsection*{Questão 23.9}
%--------------------------------------------------------------------------
\begin{questao}{Enunciado}
Explique por que os modelos de estimativa de esforço, como o COCOMO, podem
não funcionar bem quando aplicados a sistemas muito grandes.
\end{questao}
\begin{resposta}
\begin{enumerate}[nosep]
  \item \textbf{Dados históricos escassos:} COCOMO foi calibrado com projetos
        de até algumas centenas de KLOC. Sistemas com milhões de LOC são raros
        e têm características únicas não capturadas na base de dados histórica.
  \item \textbf{Não-linearidade da complexidade:} o expoente de escala ($E$)
        em $\text{PM} = A \times (\text{KLOC})^E$ assume que a complexidade
        cresce de forma regular com o tamanho. Em sistemas muito grandes, a
        complexidade de integração, testes de sistema e comunicação organizacional
        cresce de forma super-linear, além do que o expoente modela.
  \item \textbf{Divisão em subsistemas:} sistemas muito grandes são
        desenvolvidos por múltiplas equipes e organizações em paralelo,
        com interfaces complexas. O COCOMO não modela adequadamente o custo
        de integração entre grandes subsistemas.
  \item \textbf{Evolução durante o desenvolvimento:} em projetos de vários
        anos, os requisitos mudam significativamente. O LOC estimado no
        início é uma meta-alvo instável, invalidando as estimativas baseadas
        nele.
\end{enumerate}
\end{resposta}

%--------------------------------------------------------------------------
\subsection*{Questão 23.10}
%--------------------------------------------------------------------------
\begin{questao}{Enunciado}
É ético para uma empresa cotar um preço baixo para um contrato de software
sabendo que os requisitos são ambíguos e que eles poderão cobrar um preço
elevado para mudanças posteriores solicitadas pelo cliente?
\end{questao}
\begin{resposta}
\textbf{Não, esta prática é eticamente problemática} pelas seguintes razões:

\begin{itemize}[nosep]
  \item \textbf{Enganosidade intencional:} a empresa conhece a assimetria
        de informação (que os requisitos são ambíguos e gerarão mudanças
        custosas) e a explora deliberadamente. Isso viola o princípio ético
        de agir de boa-fé com o cliente.
  \item \textbf{Dano ao cliente:} o cliente toma sua decisão de compra baseado
        em um preço que não reflete o custo real. O orçamento será excedido,
        potencialmente comprometendo outros projetos ou a viabilidade
        financeira da organização compradora.
  \item \textbf{Conflito com códigos de ética profissional:} o Código de
        Ética da ACM (Princípio 1.3) determina que profissionais de computação
        devem ser honestos e não enganar clientes ou empregadores.
  \item \textbf{Efeito sistêmico:} a prática destrói a confiança no setor,
        eleva custos de contratação (clientes precisam de auditorias extensas
        e contratos mais complexos) e prejudica empresas éticas que propõem
        preços justos e perdem contratos para a concorrência desonesta.
\end{itemize}

A alternativa ética é: identificar e documentar a ambiguidade dos requisitos
antes da proposta, cotar um preço que inclua uma reserva de risco explícita,
e propor uma fase de detalhamento de requisitos paga antes da estimativa final.
\end{resposta}

%==========================================================================
\section*{Exercício 12.1 — Estimativa de Custo com COCOMO}
%==========================================================================

\begin{questao}{Enunciado}
Sistema GUI-based com 75.000 LOC (complexo). Calcule: (a) esforço total em
PM; (b) distribuição por fase; (c) esforço e custos por atividade em cada
fase; (d) custo total do projeto.
\end{questao}

\textbf{(a) Esforço total:}
\[
  \text{PM} = 3{,}0 \times (\text{KLOC})^{1{,}12}
             = 3{,}0 \times 75^{1{,}12}
             = 3{,}0 \times 125{,}91
             \approx \boxed{377{,}7 \text{ PM}}
\]

\textbf{(b) Distribuição por fase:}

\begin{center}
\begin{tabular}{lcc}
\toprule
\textbf{Fase} & \textbf{\%} & \textbf{PM} \\
\midrule
Requirements Analysis \& Conception & 8\% & 30,2 PM \\
Design & 18\% & 68,0 PM \\
Implementation & 48\% & 181,3 PM \\
Integration \& Test & 26\% & 98,2 PM \\
\midrule
\textbf{Total} & 100\% & \textbf{377,7 PM} \\
\bottomrule
\end{tabular}
\end{center}

\textbf{(c) Esforço e custos por atividade — Fase: Requirements Analysis \& Conception (30,2 PM):}

\begin{center}\small
\begin{tabular}{lcccc}
\toprule
\textbf{Atividade} & \textbf{\%} & \textbf{PM} & \textbf{Taxa (€/PM)} & \textbf{Custo (€)} \\
\midrule
Requirements Analysis  & 42\% & 12,69 & 10.000 & 126.920 \\
Product Design         & 16\% &  4,84 &  6.000 &  29.010 \\
Programming            & 10\% &  3,02 &  6.000 &  18.130 \\
Plan Tests             &  6\% &  1,81 &  9.000 &  16.320 \\
Verification \& Validation & 10\% & 3,02 & 9.000 & 27.200 \\
Project Management     &  8\% &  2,42 & 12.000 &  29.010 \\
Configuration Mgmt     &  3\% &  0,91 &  7.000 &   6.350 \\
Quality Assurance      &  5\% &  1,51 &  7.000 &  10.580 \\
\midrule
\textbf{Total RAC}     & 100\% & 30,22 & --- & \textbf{263.520} \\
\bottomrule
\end{tabular}
\end{center}

\textit{Os valores para Design (506.550 €), Implementation (1.292.770 €) e
Integration \& Test (720.870 €) foram computados analogamente usando as
colunas correspondentes da Figura~1.}

\textbf{(d) Custo total do projeto:}

\begin{center}
\begin{tabular}{lc}
\toprule
\textbf{Fase} & \textbf{Custo (€)} \\
\midrule
Requirements Analysis \& Conception & 263.520 \\
Design & 506.550 \\
Implementation & 1.292.770 \\
Integration \& Test & 720.870 \\
\midrule
\textbf{TOTAL} & \textbf{2.783.710} \\
\bottomrule
\end{tabular}
\end{center}

%==========================================================================
\section*{Exercício 12.2 — Rede de Planejamento (CPM / MPM)}
%==========================================================================

\begin{questao}{Enunciado}
Para o depósito automatizado, construa: (a) rede CPM com caminho crítico;
(b) rede MPM; (c) gráfico de barras (Gantt).
\end{questao}

\textbf{(a) Rede CPM — Análise de Caminhos Críticos}

Cálculo dos tempos mais cedo e mais tarde:

\begin{center}\small
\begin{tabular}{ccccccc}
\toprule
\textbf{Atividade} & \textbf{Dur.} & \textbf{ES} & \textbf{EF} & \textbf{LS} & \textbf{LF} & \textbf{Folga} \\
\midrule
A (Preparação)               & 2  & 0  & 2  & 0  & 2  & 0 \textbf{*} \\
B (Projeto depósito)         & 4  & 2  & 6  & 3  & 7  & 1 \\
C (Construção)               & 14 & 2  & 16 & 2  & 16 & 0 \textbf{*} \\
D (Desenv. sist. controle)   & 9  & 6  & 15 & 7  & 16 & 1 \\
E (Provisão mecânica)        & 3  & 6  & 9  & 16 & 19 & 10 \\
F (Teste componentes edif.)  & 7  & 16 & 23 & 16 & 23 & 0 \textbf{*} \\
G (Teste componentes mec.)   & 4  & 16 & 20 & 19 & 23 & 3 \\
H (Instalação edif.)         & 4  & 16 & 20 & 19 & 23 & 3 \\
I (Integração e introdução)  & 4  & 23 & 27 & 23 & 27 & 0 \textbf{*} \\
\bottomrule
\end{tabular}
\end{center}

\textbf{Caminho crítico:} $\mathbf{A \to C \to F \to I}$ \quad (duração: $2+14+7+4 = \mathbf{27}$ dias)

\textbf{Diagrama CPM (Activity-on-Node):}

\begin{center}
\begin{tikzpicture}[
  node distance=1.5cm and 2.2cm,
  act/.style={draw,rectangle,rounded corners=3pt,minimum width=2.2cm,
              minimum height=1.1cm,align=center,font=\footnotesize},
  crit/.style={act,fill=red!15,draw=red!70,very thick},
  norm/.style={act,fill=blue!8,draw=blue!50},
  arr/.style={-{Stealth[length=4pt]},thick},
  carr/.style={-{Stealth[length=4pt]},very thick,red!70}]

  \node[crit] (A) {\textbf{A}\\2d\\0|2|0|2};
  \node[norm,right=1.8cm of A] (B) {B\\4d\\2|6|3|7};
  \node[crit,below=0.6cm of B] (C) {\textbf{C}\\14d\\2|16|2|16};
  \node[norm,right=1.8cm of B] (D) {D\\9d\\6|15|7|16};
  \node[norm,below=0.6cm of D] (E) {E\\3d\\6|9|16|19};
  \node[crit,right=1.8cm of D] (F) {\textbf{F}\\7d\\16|23|16|23};
  \node[norm,below=0.6cm of F] (G) {G\\4d\\16|20|19|23};
  \node[norm,below=0.6cm of G] (H) {H\\4d\\16|20|19|23};
  \node[crit,right=1.8cm of F] (I) {\textbf{I}\\4d\\23|27|23|27};

  % Legend inside nodes: Dur d / ES|EF|LS|LF
  \draw[carr] (A) -- (C);
  \draw[arr] (A) -- (B);
  \draw[arr] (B) -- (D);
  \draw[arr] (B) -- (E);
  \draw[carr] (C) -- (F);
  \draw[arr] (C) -- (G);
  \draw[arr] (C) -- (H);
  \draw[arr] (D) -- (F);
  \draw[arr] (D) -- (G);
  \draw[arr] (D) -- (H);
  \draw[arr] (E) -- (G);
  \draw[carr] (F) -- (I);
  \draw[arr] (G) -- (I);
  \draw[arr] (H) -- (I);

  \node[below=0.15cm,font=\tiny,align=left] at (A.south)
    {ES|EF|LS|LF};
  \node[red!70,below right=0.1cm and 3cm of I,font=\small\bfseries]
    {Caminho crítico: A→C→F→I};
\end{tikzpicture}
\end{center}

\textbf{(b) Rede MPM (Metra Potential Method / Activity-on-Node):}

No MPM, os nós são as atividades e as arestas são ponderadas pela duração
do predecessor (ou pela relação temporal explícita). Cada nó mostra:
ES | EF na linha superior, LS | LF na linha inferior, com a duração ao centro.

\begin{center}
\begin{tikzpicture}[
  node distance=1.2cm and 1.8cm,
  mpm/.style={draw,rectangle,minimum width=2.0cm,minimum height=1.4cm,
              align=center,font=\tiny},
  cmpm/.style={draw,rectangle,minimum width=2.0cm,minimum height=1.4cm,
              align=center,font=\tiny,fill=red!12,draw=red!70,thick},
  arr/.style={-{Stealth[length=3.5pt]},gray!70},
  carr/.style={-{Stealth[length=3.5pt]},red!70,very thick}]

  \node[cmpm] (A) {\textbf{A}\\\textbf{dur=2}\\0 | 2\\0 | 2};
  \node[mpm,right=2.5cm of A,yshift=0.7cm] (B) {B\\dur=4\\2 | 6\\3 | 7};
  \node[cmpm,right=2.5cm of A,yshift=-0.7cm] (C) {\textbf{C}\\\textbf{dur=14}\\2 | 16\\2 | 16};
  \node[mpm,right=2.2cm of B,yshift=0.5cm] (D) {D\\dur=9\\6 | 15\\7 | 16};
  \node[mpm,right=2.2cm of B,yshift=-0.5cm] (E) {E\\dur=3\\6 | 9\\16 | 19};
  \node[cmpm,right=2.2cm of D,yshift=0.7cm] (F) {\textbf{F}\\\textbf{dur=7}\\16 | 23\\16 | 23};
  \node[mpm,right=2.2cm of D,yshift=-0.3cm] (G) {G\\dur=4\\16 | 20\\19 | 23};
  \node[mpm,right=2.2cm of E,yshift=-0.3cm] (H) {H\\dur=4\\16 | 20\\19 | 23};
  \node[cmpm,right=2.2cm of F] (I) {\textbf{I}\\\textbf{dur=4}\\23 | 27\\23 | 27};

  \draw[carr] (A) -- node[above,font=\tiny]{2} (C);
  \draw[arr]  (A) -- node[above,font=\tiny]{2} (B);
  \draw[arr]  (B) -- node[above,font=\tiny]{4} (D);
  \draw[arr]  (B) -- node[below,font=\tiny]{4} (E);
  \draw[carr] (C) -- node[above,font=\tiny]{14} (F);
  \draw[arr]  (C) -- node[above,font=\tiny]{14} (G);
  \draw[arr]  (C) -- node[below,font=\tiny]{14} (H);
  \draw[arr]  (D) -- node[above,font=\tiny]{9} (F);
  \draw[arr]  (D) -- node[font=\tiny]{9} (G);
  \draw[arr]  (D) -- node[font=\tiny]{9} (H);
  \draw[arr]  (E) -- node[font=\tiny]{3} (G);
  \draw[carr] (F) -- node[above,font=\tiny]{7} (I);
  \draw[arr]  (G) -- node[font=\tiny]{4} (I);
  \draw[arr]  (H) -- node[font=\tiny]{4} (I);
\end{tikzpicture}
\end{center}

\textbf{(c) Gráfico de Barras (Gantt) para o Depósito Automatizado:}

\begin{center}
\begin{tikzpicture}[xscale=0.43, yscale=0.48]
  \foreach \x in {0,...,27} {
    \draw[gray!25,thin] (\x,-0.3) -- (\x,9.3);
    \node[below,font=\tiny] at (\x,-0.3) {\x};
  }
  % row/name/ES/EF/LS/LF
  \foreach \row/\name/\s/\f/\ls/\lf in {
      1/A/0/2/0/2,   2/B/2/6/3/7,
      3/C/2/16/2/16, 4/D/6/15/7/16,
      5/E/6/9/16/19, 6/F/16/23/16/23,
      7/G/16/20/19/23, 8/H/16/20/19/23,
      9/I/23/27/23/27} {
    \pgfmathsetmacro{\y}{9-\row}
    % Float bar
    \fill[orange!20,draw=orange!40,dashed]
      (\f,\y+0.1) rectangle (\lf,\y+0.8);
    % Activity bar
    \fill[blue!55,draw=blue!80,rounded corners=0.5pt]
      (\s,\y+0.1) rectangle (\f,\y+0.8);
    \node[left,font=\small\bfseries] at (0,\y+0.45) {\name};
  }
  % Critical path: A(row1), C(row3), F(row6), I(row9)
  \foreach \row/\s/\f in {1/0/2, 3/2/16, 6/16/23, 9/23/27} {
    \pgfmathsetmacro{\y}{9-\row}
    \fill[red!55,draw=red!80,rounded corners=0.5pt]
      (\s,\y+0.1) rectangle (\f,\y+0.8);
  }
  \draw[->] (0,-0.3) -- (28.5,-0.3) node[right,font=\small] {Dias};
  \draw (0,-0.3) -- (0,9.3);
  \node[above,font=\small\bfseries] at (13.5,9.5) {Gantt --- Dep\'osito Automatizado};
  \fill[red!55] (15,0.5) rectangle (16,0.9);
  \node[right,font=\tiny] at (16,0.7) {Cr\'itico};
  \fill[blue!55] (15,0.0) rectangle (16,0.4);
  \node[right,font=\tiny] at (16,0.2) {Normal};
  \fill[orange!20,draw=orange!40] (19,0.5) rectangle (20,0.9);
  \node[right,font=\tiny] at (20,0.7) {Folga};
\end{tikzpicture}
\end{center}

%==========================================================================
\section*{Exercício 12.3 — Net Plans vs.\ Diagramas de Gantt}
%==========================================================================

\begin{questao}{Enunciado}
Nomeie as vantagens e desvantagens de net plans e diagramas de Gantt.
\end{questao}
\begin{resposta}
\begin{center}
\begin{tabular}{p{3cm}p{5.5cm}p{5.5cm}}
\toprule
 & \textbf{Net Plan (CPM/MPM)} & \textbf{Diagrama de Gantt} \\
\midrule
\textbf{Vantagens}
& \begin{itemize}[nosep,leftmargin=*]
    \item Representa explicitamente dependências entre atividades
    \item Identifica o caminho crítico e as folgas de cada atividade
    \item Permite análise formal do impacto de atrasos (``what-if'')
    \item Útil para projetos complexos com muitas interdependências
  \end{itemize}
& \begin{itemize}[nosep,leftmargin=*]
    \item Visualização intuitiva do cronograma temporal
    \item Fácil de ler por não-especialistas e clientes
    \item Mostra claramente quem faz o quê e quando
    \item Ferramentas amplamente disponíveis (MS Project, etc.)
  \end{itemize} \\
\midrule
\textbf{Desvantagens}
& \begin{itemize}[nosep,leftmargin=*]
    \item Difícil de ler para não-especialistas
    \item Não mostra diretamente alocação de recursos
    \item Construção e manutenção manual são trabalhosas
    \item Para projetos muito grandes, torna-se visualmente incompreensível
  \end{itemize}
& \begin{itemize}[nosep,leftmargin=*]
    \item Não representa explicitamente as dependências
    \item Não identifica o caminho crítico automaticamente
    \item Atualizações frequentes exigem redesenho
    \item Menos adequado para análise formal de riscos de prazo
  \end{itemize} \\
\bottomrule
\end{tabular}
\end{center}

\textbf{Conclusão:} Projetos complexos se beneficiam de ambas as ferramentas
complementarmente: a rede (CPM/MPM) para análise e planejamento técnico, e
o Gantt para comunicação com a equipe e stakeholders.
\end{resposta}

\vspace{8pt}\hrule\vspace{6pt}
\begin{center}
  {\small \textit{Referência:} Ian Sommerville.
  \textbf{Software Engineering}, 10ª ed. Pearson, 2016.
  Capítulos 22 (Gerenciamento de Projetos) e 23 (Planejamento de Projetos).}
\end{center}

\end{document}
